\chapter{1984年全国卷（文）}

\section{单选题}

\begin{problem}

数集 \(X = \{(2n + 1)\pi, n\in\Z\}\) 与数集 \(Y = \{(4k \pm 1)\pi, k\in\Z\}\) 之间的关系是（\(\quad\)）

\choices
    {\(X \subseteq Y\)}
    {\(X \supseteq Y\)}
    {\(X = Y\)}
    {\(X \neq Y\)}
\end{problem}

\begin{answer}
C
\end{answer}

\begin{solution}
任取 \(n\in\Z\). 当 \(n=2q\)（\(q\in\Z\)）时，\(2n+1=4q+1\)；当 \(n=2q+1\)（\(q\in\Z\)）时，\(2n+1=4(q+1)-1\)，所以 \(X\subseteq Y\). 反过来，\(4k+1=2(2k)+1\)，\(4k-1=2(2k-1)+1\)，所以 \(Y\subseteq X\). 因此 \(X=Y\)，选 C.
\end{solution}

\begin{problem}

函数 \(y = f(x)\) 与它的反函数 \(y = f^{-1}(x)\) 的图象（\(\quad\)）

\choices
    {关于 \(y\) 轴对称}
    {关于原点对称}
    {关于直线 \(x + y = 0\) 对称}
    {关于直线 \(x - y = 0\) 对称}
\end{problem}

\begin{answer}
D
\end{answer}

\begin{solution}
若点 \((x,y)\) 在函数 \(y=f(x)\) 的图象上，则点 \((y,x)\) 在反函数 \(y=f^{-1}(x)\) 的图象上. 交换点的横、纵坐标正是关于直线 \(y=x\)，即 \(x-y=0\) 的轴对称变换，因此选 D.
\end{solution}

\begin{problem}

复数 \(\frac{1}{2} - \frac{\sqrt{3}}{2}\i\) 的三角形式是（\(\quad\)）

\choices
    {\(\cos \left(-\frac{\pi}{3}\right) + \i\sin \left(-\frac{\pi}{3}\right)\)}
    {\(\cos \frac{\pi}{3} + \i\sin \frac{\pi}{3}\)}
    {\(\cos \frac{\pi}{3} - \i\sin \frac{\pi}{3}\)}
    {\(\cos \frac{\pi}{3} + \i\sin \frac{5\pi}{6}\)}
\end{problem}

\begin{answer}
A
\end{answer}

\begin{solution}
复数的模为 \(\sqrt{\left(\frac{1}{2}\right)^2+\left(-\frac{\sqrt{3}}{2}\right)^2}=1\). 其对应点的横坐标为 \(\frac{1}{2}\)，纵坐标为 \(-\frac{\sqrt{3}}{2}\)，故幅角可取 \(-\frac{\pi}{3}\). 因而 \(\frac{1}{2}-\frac{\sqrt{3}}{2}\i=\cos\left(-\frac{\pi}{3}\right)+\i\sin\left(-\frac{\pi}{3}\right)\)，按标准三角形式选 A. 由于 \(\sin\left(-\frac{\pi}{3}\right)=-\sin\frac{\pi}{3}\)，原卷的这两个选项按恒等变形也表示同一复数；这里保留原卷及官方答案 A.
\end{solution}

\begin{problem}

直线与平面平行的充要条件是这条直线与平面内的（\(\quad\)）

\choices
    {一条直线不相交}
    {两条直线不相交}
    {任意一条直线都不相交}
    {无数条直线不相交}
\end{problem}

\begin{answer}
C
\end{answer}

\begin{solution}
若直线与平面平行，则它与该平面没有公共点，从而不可能与平面内的任何一条直线相交. 反之，若一条直线与平面内任意一条直线都不相交，则它不可能在该平面内，因为平面内存在与它相交的直线；它也不可能与该平面相交，因为若交点为 \(P\)，则平面内任意一条经过 \(P\) 的直线都与它相交. 所以该直线在平面外且与平面没有公共点，即与该平面平行，故选 C.
\end{solution}

\begin{problem}

方程 \(x^{2} - 79x + 1 = 0\) 的两根可分别作为（\(\quad\)）

\choices
    {一椭圆和一双曲线的离心率}
    {两抛物线的离心率}
    {一椭圆和一抛物线的离心率}
    {两椭圆的离心率}
\end{problem}

\begin{answer}
A
\end{answer}

\begin{solution}
设方程的两根为 \(r_{1}\)，\(r_{2}\). 其判别式为 \(\Delta=79^{2}-4>0\)，故两根为实数. 由韦达定理，\(r_{1}r_{2}=1>0\)，\(r_{1}+r_{2}=79>0\)，所以两根均为正数. 由于两根的积为 \(1\)，且它们不可能同时为 \(1\)（否则和为 \(2\)），故一根大于 \(1\)，另一根小于 \(1\). 椭圆的离心率满足 \(0<e<1\)，双曲线的离心率满足 \(e>1\)，而抛物线的离心率等于 \(1\)，所以两根可分别作为一椭圆和一双曲线的离心率，选 A.
\end{solution}

\section{填空题}

\begin{problem}

已知函数 \(\log_{0.5}(2x - 3) > 0\)，求 \(x\) 的取值范围\fillinblank{}.
\end{problem}

\begin{answer}
\(\frac{3}{2}<x<2\).
\end{answer}

\begin{solution}
对数真数必须为正，所以 \(2x-3>0\). 由于底数 \(0.5\in(0,1)\)，对数函数单调递减，且 \(\log_{0.5}u>0\) 等价于 \(0<u<1\). 因此 \(0<2x-3<1\)，解得 \(\frac{3}{2}<x<2\).
\end{solution}

\begin{problem}

已知圆柱的侧面展开图是边长为 \(2\) 与 \(4\) 的矩形，求圆柱的体积\fillinblank{}.
\end{problem}

\begin{answer}
\(\frac{4}{\pi}\) 或 \(\frac{8}{\pi}\).
\end{answer}

\begin{solution}
设圆柱的底面半径为 \(r\)，高为 \(h\). 侧面展开矩形的一边是圆柱的高，另一边是底面圆的周长，因此有两种对应情形：当 \(h=2\)，\(2\pi r=4\) 时，\(r=\frac{2}{\pi}\)，从而 \(V=\pi r^{2}h=\frac{8}{\pi}\)；当 \(h=4\)，\(2\pi r=2\) 时，\(r=\frac{1}{\pi}\)，从而 \(V=\pi r^{2}h=\frac{4}{\pi}\). 所以圆柱的体积为 \(\frac{4}{\pi}\) 或 \(\frac{8}{\pi}\).
\end{solution}

\begin{problem}

已知实数 \(m\) 满足 \(2x^{2} - (2\i - 1)x + m - \i = 0\)，求 \(m\) 及 \(x\) 的值\fillinblank{}.
\end{problem}

\begin{answer}
在 \(x\in\R\) 的约定下，\(m=0\)，\(x=-\frac{1}{2}\)；若不作此约定，解不唯一.
\end{answer}

\begin{solution}
题面只给出 \(m\in\R\)；若按通常的实变量约定再取 \(x\in\R\)，将方程左边整理为 \(2x^{2}+x+m-(2x+1)\i=0\). 由于复数等于零当且仅当实部和虚部都为零，因此
\(\begin{cases}
2x^{2}+x+m=0,\\
2x+1=0.
\end{cases}\)
由第二式得 \(x=-\frac{1}{2}\)，代入第一式得 \(m=0\).

若只按题面字面条件允许 \(x\) 为复数，令 \(x=u+v\i\)，其中 \(u\)，\(v\in\R\). 比较实部和虚部得
\(\begin{cases}
2u^{2}-2v^{2}+u+2v+m=0,\\
(4u+1)v-2u-1=0.
\end{cases}\)
第二式排除 \(u=-\frac{1}{4}\)，并给出 \(v=\frac{2u+1}{4u+1}\)；再由第一式得 \(m=-2u^{2}+2v^{2}-u-2v\). 因而任取满足 \(u\in\R\) 且 \(u\ne-\frac{1}{4}\) 的 \(u\) 都得到一个解，例如 \(u=0\)，\(v=1\) 时 \(x=\i\)，\(m=0\). 所以 \(m=0\)，\(x=-\frac{1}{2}\) 只有在补充 \(x\in\R\) 的约定下才是唯一答案.
\end{solution}

\begin{problem}

求 \(\displaystyle\lim_{n\to\infty}\frac{(n^2 + 1) + (n^2 + 2) + \cdots + (n^2 + n)}{n(n - 1)(n - 2)}\) 的值\fillinblank{}.
\end{problem}

\begin{answer}
\(1\).
\end{answer}

\begin{solution}
分子为 \((n^{2}+1)+(n^{2}+2)+\cdots+(n^{2}+n)=n^{3}+\frac{n(n+1)}{2}\)，分母为 \(n(n-1)(n-2)=n^{3}-3n^{2}+2n\). 因而 \(\displaystyle\frac{n^{3}+\frac{n(n+1)}{2}}{n^{3}-3n^{2}+2n}=\frac{1+\frac{n+1}{2n^{2}}}{1-\frac{3}{n}+\frac{2}{n^{2}}}\to1\)，所求极限为 \(1\).
\end{solution}

\begin{problem}

要排一张有 \(6\) 个歌唱节目和 \(4\) 个舞蹈节目的演出节目单，任何两个舞蹈节目不得相邻，问有多少种不同的排法\fillinblank{}.
\end{problem}

\begin{answer}
\(\mathrm A_7^4\cdot 6!\).
\end{answer}

\begin{solution}
先排定 \(6\) 个不同的歌唱节目，有 \(6!\) 种排法. 歌唱节目排好后形成 \(7\) 个空位，包括两端的空位；为使舞蹈节目不相邻，每个空位至多放一个舞蹈节目. 依次把 \(4\) 个不同的舞蹈节目放入 \(7\) 个空位，有 \(\mathrm A_7^4\) 种放法. 因而不同的节目单共有 \(\mathrm A_7^4\cdot6!\) 种.
\end{solution}

\section{解答题}

\begin{problem}

画出方程 \(y^{2} = -4x \) 的曲线.
\end{problem}

\begin{solution}
方程可写成 \(y^{2}=4px\)，其中 \(p=-1\). 因此该曲线是顶点为 \(O(0,0)\)、焦点为 \((-1,0)\)、准线为 \(x=1\) 的抛物线，开口向左，并且关于 \(x\) 轴对称；其图象经过 \((-1,2)\)、\((-1,-2)\) 等点. 按上述顶点、对称轴和开口方向作图即可.
\end{solution}

\begin{problem}

画出函数 \(y = \frac{1}{(x + 1)^2} \) 的图象.
\end{problem}

\begin{solution}
函数的定义域为 \(x\ne-1\)，且 \(y>0\). 直线 \(x=-1\) 与 \(y=0\) 分别是它的竖直渐近线和水平渐近线；图象关于直线 \(x=-1\) 对称，并经过 \((-2,1)\) 与 \((0,1)\). 又 \(\displaystyle\frac{\mathrm{d}y}{\mathrm{d}x}=-\frac{2}{(x+1)^{3}}\)，所以函数在 \((-\infty,-1)\) 上递增，在 \((-1,+\infty)\) 上递减. 因此图象由位于 \(x\) 轴上方、分别在直线 \(x=-1\) 两侧的两支曲线组成.
\end{solution}

\begin{problem}

已知等差数列 \(a\)，\(b\)，\(c\) 中的三个数都是正数，且公差不为零. 求证它们的倒数所组成的数列 \(\frac{1}{a}\)，\(\frac{1}{b}\)，\(\frac{1}{c}\) 不可能成等差数列.
\end{problem}

\begin{solution}
由 \(a\)，\(b\)，\(c\) 成等差数列，得 \(a+c=2b\). 又因公差不为零，所以 \(a\ne c\). 假设 \(\frac{1}{a}\)，\(\frac{1}{b}\)，\(\frac{1}{c}\) 也成等差数列，则 \(\displaystyle\frac{2}{b}=\frac{1}{a}+\frac{1}{c}=\frac{a+c}{ac}=\frac{2b}{ac}\). 由于 \(a\)，\(b\)，\(c\) 均为正数，约去 \(2b\) 得 \(ac=b^{2}\). 联立 \(a+c=2b\)，有 \(\displaystyle(a-c)^{2}=(a+c)^{2}-4ac=4b^{2}-4b^{2}=0\)，从而 \(a=c\)，这与公差不为零矛盾. 所以三数的倒数不可能成等差数列.
\end{solution}

\begin{problem}

把 \(1 - \frac{1}{4}\sin^2 2\alpha -\sin^2\beta -\cos^4\alpha\) 化成三角函数的积的形式（要求结果最简）.
\end{problem}

\begin{solution}
原式 \(=1-\sin^{2}\beta-\sin^{2}\alpha\cos^{2}\alpha-\cos^{4}\alpha\)
\(=\cos^{2}\beta-\cos^{2}\alpha(\sin^{2}\alpha+\cos^{2}\alpha)=\cos^{2}\beta-\cos^{2}\alpha\)
\(=(\cos\beta-\cos\alpha)(\cos\beta+\cos\alpha)\)
\(=\left(2\sin\frac{\alpha+\beta}{2}\sin\frac{\alpha-\beta}{2}\right)\left(2\cos\frac{\alpha+\beta}{2}\cos\frac{\alpha-\beta}{2}\right)\)
\(=\sin(\alpha+\beta)\sin(\alpha-\beta)\).
\end{solution}

\begin{problem}

如图，经过正三棱柱底面一边 \(AB\)，作与底面成 \(30^{\circ}\) 角的平面，已知截面三角形 \(ABD\) 的面积为 \(32\mathrm{~cm}^{2}\)，求截得的三棱锥 \(D - ABC\) 的体积.

\bitmapfigure[max width=0.34\linewidth]{img/1984/national_liberal/q15_fig1.png}
\end{problem}

\begin{solution}
在底面三角形 \(ABC\) 中作 \(CE\perp AB\)，并连接 \(DE\). 由于三棱柱是正三棱柱，\(ABC\) 为正三角形，且 \(DC\perp\) 底面；由三垂线定理，\(DE\perp AB\). 因此截面与底面的二面角的平面角为 \(\angle DEC=30^{\circ}\).

设 \(AB=a\). 由 \(ABC\) 为正三角形，\(CE=\frac{\sqrt{3}}{2}a\). 在直角三角形 \(DCE\) 中，\(\displaystyle DE=\frac{CE}{\cos30^{\circ}}=a\)，且 \(DC=DE\sin30^{\circ}=\frac{a}{2}\). 于是 \(\displaystyle32=S_{\triangle ABD}=\frac{1}{2}AB\cdot DE=\frac{a^{2}}{2}\)，得 \(a=8\mathrm{~cm}\).

所以底面面积为 \(S_{\triangle ABC}=\frac{\sqrt{3}}{4}a^{2}=16\sqrt{3}\mathrm{~cm}^{2}\)，三棱锥的高为 \(DC=4\mathrm{~cm}\). 因而 \(\displaystyle V=\frac{1}{3}S_{\triangle ABC}\cdot DC=\frac{1}{3}\cdot16\sqrt{3}\cdot4=\frac{64\sqrt{3}}{3}\mathrm{~cm}^{3}\).
\end{solution}

\begin{problem}

某工厂 \(1983\) 年生产某种产品 \(2\text{ 万件}\)，计划从 \(1984\) 年开始，每年的产量比上一年增长 \(20\%\). 问从哪一年开始，这家工厂生产这种产品的年产量超过 \(12\text{ 万件}\)（已知 \(\lg 2 = 0.3010\)，\(\lg 3 = 0.4771\)）.
\end{problem}

\begin{solution}
以 \(1983\) 年为第 \(1\) 年，设第 \(n\) 年的产量为 \(a_n\)（单位：万件）. 则 \(a_1=2\)，且 \(a_n=2(1.2)^{n-1}\). 令产量等于 \(12\) 万件，得 \(\displaystyle2(1.2)^{n-1}=12\)，即 \(n=1+\frac{\lg6}{\lg1.2}\).

由已知数据，\(\lg6=\lg2+\lg3=0.7781\)，\(\lg5=1-\lg2=0.6990\)，所以 \(\lg1.2=\lg\frac{6}{5}=0.0791\). 因而 \(\displaystyle n=1+\frac{0.7781}{0.0791}\approx10.84\).

数列 \(\{a_n\}\) 的公比 \(1.2>1\)，所以它递增. 产量超过 \(12\) 万件的条件是 \(n>10.84\)，而 \(n\) 为正整数，故最小的 \(n\) 是 \(11\). 第 \(11\) 年对应 \(1983+10=1993\) 年，故从 \(1993\) 年开始.
\end{solution}

\begin{problem}

已知两个椭圆的方程分别是 \(C_{1}\)：\(x^{2} + 9y^{2} - 45 = 0\)，\(C_{2}\)：\(x^{2} + 9y^{2} - 6x - 27 = 0\).

\begin{enumerate}
\item 求这两个椭圆的中心，焦点的坐标；
\item 求经过这两个椭圆的交点且与直线 \(x - 2y + 11 = 0\) 相切的圆的方程.
\end{enumerate}
\end{problem}

\begin{solution}
\begin{enumerate}
\item \(C_{1}\) 化为 \(\displaystyle\frac{x^{2}}{45}+\frac{y^{2}}{5}=1\)，故中心为 \((0,0)\)，焦半距为 \(\sqrt{45-5}=2\sqrt{10}\)，焦点为 \((2\sqrt{10},0)\)，\((-2\sqrt{10},0)\). \(C_{2}\) 化为 \(\displaystyle\frac{(x-3)^{2}}{36}+\frac{y^{2}}{4}=1\)，故中心为 \((3,0)\)，焦半距为 \(\sqrt{36-4}=4\sqrt{2}\)，焦点为 \((3+4\sqrt{2},0)\)，\((3-4\sqrt{2},0)\).
\item 两椭圆的交点满足 \(45=6x+27\)，所以 \(x=3\)，再代入 \(x^{2}+9y^{2}=45\) 得 \(y=\pm2\). 设所求圆为 \(x^{2}+y^{2}+Dx+Ey+F=0\)，代入 \((3,2)\) 与 \((3,-2)\)，分别得 \(13+3D+2E+F=0\)，\(13+3D-2E+F=0\)，相减得 \(E=0\)，再得 \(F=-3D-13\). 所以圆族为 \(x^{2}+y^{2}+Dx-3D-13=0\).

将直线 \(x-2y+11=0\) 写为 \(y=\frac{x+11}{2}\)，代入圆族并乘以 \(4\)，得 \(\displaystyle5x^{2}+(22+4D)x+69-12D=0\). 相切时该关于 \(x\) 的二次方程有重根，故 \(\displaystyle(22+4D)^{2}-20(69-12D)=0\)，即 \(D^{2}+26D-56=0\). 解得 \(D=2\) 或 \(D=-28\). 因而所求圆为 \(x^{2}+y^{2}+2x-19=0\) 或 \(x^{2}+y^{2}-28x+71=0\).
\end{enumerate}
\end{solution}
